\documentclass[11pt]{article}
\usepackage[margin=1in]{geometry}
\usepackage{amsmath,amssymb,amsthm,mathtools,hyperref}
\newtheorem{theorem}{Theorem}
\newtheorem{lemma}{Lemma}
\newcommand{\Z}{\mathbb Z}
\newcommand{\Q}{\mathbb Q}
\title{A denominator bound for a contiguous basic hypergeometric ratio}
\author{Principia Math In Progress}
\date{27 September 2026}
\begin{document}
\maketitle
\begin{abstract}
We prove Kim and Stanton's Conjecture 8.3 in the final 2023 version of
\emph{Three families of $q$-Lommel polynomials}. A quadratic coefficient
recurrence, obtained directly from the basic hypergeometric equation,
preserves the conjectured denominator bound. The argument is over an
integer polynomial ring and applies to every coefficient.
\end{abstract}

\section*{Task 1: the conjecture and its resolution}
Let $a,b,c,q$ be indeterminates, put $K=\Q(a,b,c,q)$, and work in $K[[z]]$.
Write $(x;q)_m=\prod_{j=0}^{m-1}(1-xq^j)$, with empty product $1$, and set
\[
 F(z)={}_2\phi_1(a,b;c;q,z)
 =\sum_{m\geq0}\frac{(a;q)_m(b;q)_m}{(c;q)_m(q;q)_m}z^m,
 \qquad G(z)={}_2\phi_1(aq,bq;cq;q,z).
\]
Since $F(0)=1$, its inverse exists. Define
\[
 R(z)=\frac{G(z)}{(1-c)F(z)}=\sum_{n\geq0}r_nz^n,
 \qquad
 D_n=\prod_{k=0}^n(1-cq^k)^{\lfloor(n+1)/(k+1)\rfloor}.
\]
Thus $r_n=\gamma_n(a,b,c)/(1-c)$ in the notation of the conjecture.

\begin{theorem}\label{thm:main}
For every $n\geq0$, $P_n:=D_nr_n$ belongs to $\Z[a,b,c,q]$.
Consequently MathDB 350746, equivalently Conjecture 8.3 of
Kim--Stanton~\cite{KS}, holds in full.
\end{theorem}

The assertion is a denominator bound, not an assertion that $D_n$ is
the reduced denominator or that $P_n$ has nonnegative coefficients.
The latter is not part of the source conjecture.

\begin{lemma}\label{lem:rec}
Put $A=(1-a)(1-b)$. Then $r_0=1/(1-c)$ and, for $n\geq1$,
\begin{align}
 (1-cq^n)r_n={}&\bigl(a+b-ab-abq^n\bigr)r_{n-1}
 -Ac\sum_{i+j=n-1}q^j r_i r_j\notag\\
 &+Aabq\sum_{i+j=n-2}q^j r_i r_j .\label{eq:rec}
\end{align}
Every summation index is nonnegative; a sum with negative total is zero.
\end{lemma}
\begin{proof}
If $f_m=[z^m]F(z)$, the defining product gives, for $m\geq1$,
\[
 (1-q^m)(1-cq^{m-1})f_m
 =(1-aq^{m-1})(1-bq^{m-1})f_{m-1}.
\]
Multiplying this equality by $qz^m$ and summing, or simply comparing
coefficients, yields
\begin{equation}
 (c-abqz)F(q^2z)-\bigl(c+q-(a+b)qz\bigr)F(qz)
       +q(1-z)F(z)=0.\label{eq:linear}
\end{equation}
The same product formula also gives
\begin{equation}
 F(z)-F(qz)=\frac{Az}{1-c}G(z).\label{eq:shift}
\end{equation}
In fact, the coefficient of $z^{m+1}$ on the left is
$(1-q^{m+1})f_{m+1}=A[z^m]G(z)/(1-c)$.

Let $H(z)=F(qz)/F(z)$. Equation~\eqref{eq:shift} says
$H(z)=1-AzR(z)$, and $F(q^2z)/F(z)=H(qz)H(z)$.
Divide~\eqref{eq:linear} by $F(z)$, substitute these expressions,
and cancel the common factor $Aqz$. This is legitimate over $K[[z]]$,
where $Aq$ is nonzero and multiplication by $z$ is injective. The result is
\begin{align}
 R(z)-cR(qz)={}&1+z\bigl((a+b-ab)R(z)-abqR(qz)\bigr)\notag\\
 &-AczR(z)R(qz)+Aabqz^2R(z)R(qz).\label{eq:riccati}
\end{align}
The constant coefficient gives $r_0=1/(1-c)$. Taking the coefficient of
$z^n$ for $n\geq1$ gives~\eqref{eq:rec}.
\end{proof}

\begin{lemma}\label{lem:den}
For $n\geq1$, put $E_n=D_n/(1-cq^n)$. In $\Z[c,q]$ one has
\[
 D_{n-1}\mid E_n,
 \qquad D_iD_j\mid E_n\quad\text{if }i,j\geq0\text{ and }i+j\leq n-1.
\]
\end{lemma}
\begin{proof}
The factor with $k=n$ in $D_n$ has exponent $1$, so
$E_n=\prod_{k=0}^{n-1}(1-cq^k)^{\lfloor(n+1)/(k+1)\rfloor}$.
For every $0\leq k<n$,
\[
 \left\lfloor\frac{n}{k+1}\right\rfloor
 \leq\left\lfloor\frac{n+1}{k+1}\right\rfloor,
\]
which proves the first divisibility. The exponent of the displayed factor
in the product presentation of $D_iD_j$ is
\[
 \left\lfloor\frac{i+1}{k+1}\right\rfloor+
 \left\lfloor\frac{j+1}{k+1}\right\rfloor
 \leq\left\lfloor\frac{i+j+2}{k+1}\right\rfloor
 \leq\left\lfloor\frac{n+1}{k+1}\right\rfloor.
\]
Here a floor is zero when $k$ exceeds the relevant index. As $i,j<n$,
neither product has a factor indexed by $k\geq n$. Thus each quotient is
explicitly a product of factors $1-cq^k$ with nonnegative integer exponents.
No coprimality assertion is needed.
\end{proof}

\begin{proof}[Proof of Theorem~\ref{thm:main}]
We have $D_0r_0=1$. Suppose $P_i=D_ir_i\in\Z[a,b,c,q]$ for all $i<n$.
Multiplication of~\eqref{eq:rec} by $E_n$ gives the constructive recurrence
\begin{align}
 P_n={}&(a+b-ab-abq^n)\frac{E_n}{D_{n-1}}P_{n-1}\notag\\
 &-Ac\sum_{i+j=n-1}q^j\frac{E_n}{D_iD_j}P_iP_j\notag\\
 &+Aabq\sum_{i+j=n-2}q^j\frac{E_n}{D_iD_j}P_iP_j.\label{eq:P}
\end{align}
Every quotient in~\eqref{eq:P} lies in $\Z[c,q]$ by
Lemma~\ref{lem:den}. All remaining factors are integer polynomials by
induction. This proves $P_n\in\Z[a,b,c,q]$ and completes the induction.
\end{proof}

For example,
\[
 P_0=1,\qquad P_1=a+b-ab(1+q)-c+abcq.
\]
The proof is a formal rational-function identity. It therefore covers
all specializations where the coefficients are defined, including $a=1$
or $b=1$; the temporary generic cancellation of $A$ loses no such cases.
At poles, the assertion retains its rational-function meaning. Values
such as $q=1$, where the original coefficient formula is singular, may
only be interpreted after any indicated rational continuation.

\paragraph{Relation to earlier work.}
The recurrence method extends the $q$-Kishore strategy used by Kim and
Stanton~\cite{KSB}, especially their Proposition 3.3 and Theorem 3.2 for
Jackson's first $q$-Bessel function. Their earlier manuscript states the
present full-parameter assertion as Conjecture 7.2; the final paper
\cite{KS} retains it as Conjecture 8.3. The introduction of the final
paper attributes the earlier $q$-Kishore results to Lalanne's work of
1992--1993; those special cases are not new results here. The new step is
\eqref{eq:rec} for arbitrary $a,b,c$, with the integral induction
\eqref{eq:P}. All identities used in the proof have been derived above.

\section*{Proposed follow-up Tasks 2--5}
These are research directions, not additional solved claims.
\begin{enumerate}
\setcounter{enumi}{1}
\item Determine, for each $n$ and $0\leq k\leq n$, the exact generic pole
order of $r_n$ along $1-cq^k=0$, and hence the reduced denominator.
\item Find a finite weighted-tree formula for $P_n$ realizing all three
terms of~\eqref{eq:P}, with polynomial edge or vertex weights.
\item Classify the specializations $a=q^{-u}, b=q^{-v}$, for nonnegative
integers $u,v$, for which the pole orders are smaller than their generic
values, and give the resulting reduced denominator in $\Q(q)[c]$.
\item For fixed nonnegative integers $u,v,w$, determine a denominator
bound in $\Z[a,b,c,q]$ for the coefficients of
${}_2\phi_1(aq^u,bq^v;cq^w;q,z)/{}_2\phi_1(a,b;c;q,z)$,
including necessary factors that depend only on $q$.
\end{enumerate}

\begin{thebibliography}{9}
\bibitem{KS} J. S. Kim and D. Stanton,
\emph{Three families of $q$-Lommel polynomials},
S\'eminaire Lotharingien de Combinatoire \textbf{87B} (2023), Article 8,
25 pp., Conjecture 8.3, p.~24.
\href{https://www.mat.univie.ac.at/~slc/wpapers/lattpath/KimStanton/KimStanton.pdf}{Final article}.
\bibitem{KSB} J. S. Kim and D. Stanton,
\emph{Ratios of Jackson's $q$-Bessel functions and $q$-Lommel polynomials},
author manuscript dated 7 September 2020, 25 pp., Sections 2--3 and
Conjecture 7.2, pp.~23--24.
\href{https://www-users.cse.umn.edu/~stant001/PAPERS/bessel_functions_final.pdf}{Author manuscript}.
\end{thebibliography}
\end{document}
